\documentclass[12pt]{article}

% Layout.
\usepackage[top=1in, bottom=0.75in, left=1in, right=1in, headheight=1in, headsep=6pt]{geometry}

% Fonts.
\usepackage{mathptmx}
\usepackage[scaled=0.86]{helvet}
\renewcommand{\emph}[1]{\textsf{\textbf{#1}}}

% TiKZ.
\usepackage{tikz, pgfplots}
\usetikzlibrary{calc}

% Misc packages.
\usepackage{amsmath,amssymb,latexsym}
\usepackage{graphicx}
\usepackage{array}
\usepackage{xcolor}
\usepackage{multicol}
\usepackage{fancyhdr}
\pagestyle{fancy}

\usepackage[parfill]{parskip}


% Misc.
\renewcommand{\d}{\displaystyle}
\newcommand{\ds}{\displaystyle}
\def\bc{\begin{center}}
\def\ec{\end{center}}
\def\be{\begin{enumerate}}
\def\ee{\end{enumerate}}

\newcommand{\ans}[1][1in]{\rule{#1}{.5pt}}


\newcommand{\corrections}[3][l]{%
  % optional #1 = alignment: l (default), r, c
  % #2 = height (with units, e.g. 1in)
  % #3 = proportion of \linewidth (e.g. .6)
  \par\noindent
  \makebox[\linewidth][#1]{%
    \begin{tikzpicture}[baseline]
      \draw (0,0) rectangle (#3\linewidth,-#2);
      \node[anchor=south west, font=\it\small, inner sep = 0] at (0,1pt) {Write corrections here};
    \end{tikzpicture}%
  }%
  \par\vspace{2pt}
}


\lhead{Math F113X: Quiz 1}

\begin{document}

\vspace{1cm}

\noindent \textbf{Name:} \hrulefill \quad  score:\ans[1cm] \ / 10 \\


There are 10 points possible on this quiz. No aids (book, notes, etc.)
are permitted. You may use a non-programmable calculator.  {\bf Show
all work and supporting calculations for full credit. Explain how you
get your answers.}

Do not write in the ``write corrections here'' boxes while you're taking the main quiz.

\newcommand{\prefSch}{
\begin{tabular}{|c|| c|c|c|c|c|}
\hline
number of voters& 40&30&20&5&5\\
\hline \hline
1st choice&B&C&D&D&B\\ \hline
2nd choice&A&A&A&C&C\\ \hline
3rd choice&D&D&B&A&A\\ \hline
4th choice&C&B&C&B&D\\ \hline
\end{tabular}
}


\be
\item (4 points) A state is holding elections for governor. There are
  four candidates (A, B, C, and D for convenience). The preference
  schedule is below.
  
  \prefSch

\def\y{1.2}

For each of the following, provide supporting calculations.

\be
\item (1 point) How many voters voted in this election? \ans
% \vspace{1cm}
  
  \corrections[r]{1cm}{.5}
	
% \item How many voters are needed for a majority? \ans
%   \vfill
	
% \item What is the smallest number of voters that could give a
%   candidate a plurality? \ans
%   \vfill
  
\item (1 point) Find a winner under the plurality method. Show some work. \ans
%  \vspace{1cm}

  \corrections[r]{1.5cm}{.5}


%\newpage
\item (2 points) % Here's the preference schedule again.
%
%\prefSch
%
 Find a winner under Instant Runoff Voting. \ans \\
  Show some work for each round and clearly state who gets eliminated.

 
\corrections[r]{3in}{.5}

	\ee
\newpage	
\item (6 points) Below is the same preference schedule. 

\prefSch

%\begin{tabular}{|c|| c|c|c|c|c|}
%\hline
%number of voters& 40&30&20&5&5\\
%\hline \hline
%1st choice&B&C&D&D&B\\ \hline
%2nd choice&A&A&A&C&C\\ \hline
%3rd choice&D&D&B&A&A\\ \hline
%4th choice&C&B&C&B&D\\ \hline
%\end{tabular}
%\def\y{1.2}

\be
\item (2 points) Find a winner under the Borda Count method. Show some work. \ans

\vspace{1in}

\corrections{1in}{1}
  
\item (2 point) In a one-to-one comparison, who is preferred, candidate A or
  candidate B? (You must show your calculation.)
	
%  \vfill
	
  Candidate \ans\ is preferred.

  \corrections[r]{1in}{.5}
  
  
% \item Explain why candidate B cannot be the Condorcet winner.
%   \vspace{1.5in}

\item (2 points) Fill in the your result from part (b) and determine the winner under Copeland's method.
  The results of each one-on-one
  comparison (except A vs B) are provided below.
  
  The winner under Copeland's Method is \ans

  \begin{tabular}{|c||c|c|c|c|c|c|}
    \hline
    matchup & A vs B & A vs C & A vs D & B vs C & B vs D & C vs D \\ \hline
    tally & A:\hspace{0.6cm} & A: 60 & A: 75 & B: 65 & B: 45 & C: 35 \\
            & B:\hspace{0.6cm} & C: 40 & D: 25 & C: 35 & D: 55 & D: 65 \\ \hline
    winner & & A & A & B & D & D \\ \hline 
    \end{tabular}
    
	 \corrections[r]{1.1in}{.5}
	\ee	
\ee
\end{document}
%%% Local Variables:
%%% mode: LaTeX
%%% TeX-master: t
%%% End:
